\documentclass[12pt,a4paper]{article}
\usepackage[utf8]{inputenc}
\usepackage[T1]{fontenc}
\usepackage{lmodern}
\usepackage{amsmath,amssymb}
\usepackage{siunitx}
\usepackage{geometry}
\usepackage{hyperref}
\usepackage{enumitem}
\geometry{margin=2.2cm}

\title{Lecture 4 -- Offset and Bias Error Analysis}
\author{Readable Assignment Set}
\date{}

\begin{document}
\maketitle
\tableofcontents

\section{Assignments}
\subsection{D1}
Find worst-case DC output voltage of the inverting amplifier in Fig. 2.34a for $v_{in}=0$, with maximum bias current \SI{200}{nA}, maximum offset current magnitude \SI{50}{nA}, and maximum offset voltage magnitude \SI{4}{mV}.

\subsection{D2}
For AC-coupled amplifier in Fig. P2.59: explain why it is a poor AC-coupling approach (hint: bias current). Show one added component (with value) to remove bias-current effect on output.

\subsection{D3}
Given $R_1=\SI{10}{k\Omega}$ and $R_2=\SI{100}{k\Omega}$ with desired $|V_{o,DC}|\le\SI{100}{mV}$ for zero source DC input:
\begin{enumerate}[label=\alph*)]
\item max allowable op-amp offset voltage,
\item max allowable bias current,
\item add resistor to cancel bias-current effects,
\item with that resistor and ignoring $V_{OS}$, max allowable offset current.
\end{enumerate}

\subsection{D4}
Non-inverting amplifier with \(\mu\)A741C at \SI{25}{\celsius}, $\pm$\SI{15}{V}, $R_1=\SI{10}{k\Omega}$ and $R_2=\SI{470}{k\Omega}$. Ignore source resistance and finite gain/bandwidth, but include offset and bias errors. Compute maximum output error.

\subsection{D5}
Inverting amplifier with \(\mu\)A741C, same environment and values $R_1=\SI{10}{k\Omega}$, $R_2=\SI{470}{k\Omega}$.
\begin{enumerate}
\item Max output error from offset and bias.
\item Explain why adding resistor $R_3$ at non-inverting input helps.
\item Choose $R_3$ for minimum output error.
\item Find new maximum output error.
\end{enumerate}

\subsection{D6}
Inverting amplifier with $R_1=\SI{10}{k\Omega}$, $R_2=\SI{20}{k\Omega}$, $R_L=\SI{10}{k\Omega}$ and non-ideal open-loop:
$A_{0,OL}=100\,\text{dB}$, $f_{BOL}=\SI{10}{Hz}$.
Tasks: find $\alpha$, $\beta$, $\alpha/\beta$; write $20\log|A_{CL}|$; estimate $K_f$ for large/small loop gain; sketch Bode magnitude for requested terms.

\section{Hints}
\begin{itemize}
\item Worst case means signs of errors are chosen to maximize output magnitude.
\item Output error from input offset voltage scales with noise gain $1+R_2/R_1$ (non-inverting style) or equivalent closed-loop factor.
\item Bias-current error is found by multiplying bias currents with Thevenin resistances seen at op-amp input terminals.
\item To cancel average bias-current effect, make resistance seen by $v_+$ equal to resistance seen by $v_-$.
\item For D6, use $A_{CL}=\dfrac{\alpha A_{OL}}{1+\beta A_{OL}}=\dfrac{\alpha}{\beta}K_f$.
\end{itemize}

\section{Step-by-step solutions}
\subsection{D1}
\begin{enumerate}
\item Compute output contribution from $V_{OS}$ with noise gain.
\item Compute output contribution from average bias current and from offset current through feedback/input resistive network.
\item Sum magnitudes in worst-case sign combination.
\end{enumerate}

\subsection{D2}
\begin{enumerate}
\item In poor AC-coupled design, input capacitor blocks source DC but bias current has no proper DC return path, forcing node drift/saturation.
\item Add a resistor from op-amp input node to ground (or bias reference) so input bias current has a defined path.
\item Choose this resistor to match equivalent resistance at opposite input to reduce output offset.
\end{enumerate}

\subsection{D3}
Given gain magnitude $R_2/R_1=10$ and noise gain $1+R_2/R_1=11$:
\begin{enumerate}[label=\alph*)]
\item $|V_{o,OS}|=(1+R_2/R_1)|V_{OS}|\le\SI{0.1}{V}$
\[
|V_{OS}|_{max}=0.1/11=\SI{9.09}{mV}.
\]
\item Ignoring other errors, $|V_{o,IB}|\approx I_B R_2\le\SI{0.1}{V}$ gives
\[
I_{B,max}=0.1/100\,\text{k}\Omega=\SI{1}{\micro A}.
\]
\item Add resistor at non-inverting input:
\[
R_3=R_1\parallel R_2=\frac{10k\cdot100k}{110k}=\SI{9.09}{k\Omega}.
\]
\item Then only offset current remains:
\[
|V_{o,IOS}|\approx I_{OS}R_2\le\SI{0.1}{V}\Rightarrow I_{OS,max}=\SI{1}{\micro A}.
\]
\end{enumerate}

\subsection{D4}
For non-inverting gain:
\[
A_{CL}=1+\frac{R_2}{R_1}=1+\frac{470k}{10k}=48.
\]
Worst-case output offset from \(\mu\)A741C imperfections:
\[
|V_{o,max}|\approx A_{CL}|V_{OS,max}| + R_{eq,out}\cdot |I_{error,max}|,
\]
where $R_{eq,out}$ term comes from bias-current induced differential input voltage through input network. Insert datasheet worst-case values at \SI{25}{\celsius}.

\subsection{D5}
\begin{enumerate}
\item Without $R_3$, compute output offset from $V_{OS}$ plus bias/offset-current drop through effective input resistances.
\item Idea: make equal DC resistance seen by both op-amp inputs so equal bias currents produce nearly equal drops and cancel at differential input.
\item Optimal resistor:
\[
R_3=R_1\parallel R_2=\frac{10k\cdot470k}{480k}=\SI{9.79}{k\Omega}.
\]
\item Recompute with only mismatch current ($I_{OS}$) and $V_{OS}$ terms dominating. Total worst-case is significantly reduced.
\end{enumerate}

\subsection{D6}
\begin{enumerate}
\item Determine ideal closed-loop gain: $\alpha/\beta\approx -R_2/R_1=-2$ for inverting connection.
\item Feedback factor magnitude for inverting network (source grounded):
\[
\beta=\frac{R_1}{R_1+R_2}=\frac{1}{3}.
\]
\item Write
\[
20\log|A_{CL}|=20\log\left|\frac{\alpha}{\beta}\right|+20\log|K_f|.
\]
\item For $\beta A_{OL}\gg1$, $K_f\approx1$ (0 dB). For $\beta A_{OL}\ll1$, $K_f\approx\beta A_{OL}$.
\item Sketch in order: $A_{OL}$, $\alpha$, $1/\beta$, $\alpha/\beta$, $\beta A_{OL}$, $K_f$, $A_{CL}$.
\end{enumerate}

\end{document}
